% !TEX root = ../Thesis.tex

%% Leitfragen für das Abstract:
%% - Motivation: Warum ist das Thema relevant?
%% - Fragestellung: Was untersucht die Arbeit?
%% - Methode: Wie wurde untersucht (Ansatz, Techniken)?
%% - Ergebnisse: Zu welchen Schlussfolgerungen kam die Forschung?
%% - Relevanz: Was trägt die Arbeit zum bestehenden Wissen bei?
%%
%% Format: 150-250 Wörter, prägnant und eigenständig verständlich

\chapter*{Kurzfassung}
\addcontentsline{toc}{chapter}{Kurzfassung}

% Reduziere den Absatzabstand lokal, damit alles auf eine Seite passt
\begingroup
\setlength{\parskip}{0.5em}
\textbf{Kontext und Motivation:} Large Language Models (LLMs) entwickeln sich zunehmend von reinen Textgeneratoren zu agentischen Systemen, die planen, Werkzeuge einbinden und mehrschrittige Aufgaben bearbeiten können. Für das Software Engineering eröffnet dies neue Möglichkeiten zur Unterstützung von Code-Review, Testgenerierung, Fehlersuche und Refactoring.

\textbf{Zielsetzung:} Die Arbeit untersucht die Entwicklung und Evaluation agentischer Architekturen für Software-Engineering-Aufgaben. Im Zentrum stehen drei Fragen: (1) Wie lassen sich robuste Agentenstrategien für SE-Workflows systematisch modellieren? (2) Welche Architekturprinzipien ermöglichen eine sichere und nachvollziehbare Werkzeugintegration? (3) Mit welchen Metriken kann die Leistungsfähigkeit agentischer Systeme unter realistischen Randbedingungen bewertet werden?

\textbf{Methodik:} Auf Basis einer strukturierten Literaturanalyse wird eine Referenzarchitektur entwickelt, die Planung, Werkzeugnutzung, episodisches Gedächtnis und Sicherheitsmechanismen integriert. Darauf aufbauend wird ein prototypisches System in Python implementiert und anhand reproduzierbarer Szenarien aus dem Software Engineering evaluiert.

\textbf{Ergebnisse:} Die Evaluation zeigt für das gewählte Szenarienset, dass die entwickelte Architektur bei automatisiertem Refactoring eine Erfolgsrate von \pct{73} erreicht und die mittlere Bearbeitungszeit um \pct{45} reduziert. Reflexionsmechanismen verbessern die Robustheit gegenüber fehlerhaften Werkzeugausgaben um \pct{28}; optimierte Kontextverwaltung senkt die Token-Nutzung um bis zu \pct{40}. Die Resultate sind als indikative Ergebnisse einer prototypischen Umsetzung zu interpretieren und nicht als unmittelbar auf öffentliche Leaderboards übertragbare Vergleichswerte.

\textbf{Beitrag:} Die Arbeit liefert eine praxisnahe Referenzarchitektur, konkrete Implementierungsrichtlinien und ein mehrdimensionales Evaluationsschema für agentische SE-Systeme. Sie zeigt technische Potenziale, benennt methodische Grenzen und ordnet die Ergebnisse in die aktuelle Forschung zu agentischer Softwareentwicklung ein.%
\par
\endgroup

% Englisches Abstract (Standard in der Informatik)
\chapter*{Abstract}
\addcontentsline{toc}{chapter}{Abstract}

\begingroup
\setlength{\parskip}{0.5em}
\textbf{Context and Motivation:} Large Language Models (LLMs) are increasingly evolving from pure text generators into agentic systems that can plan, invoke tools, and handle multi-step tasks. In software engineering, this creates new opportunities to support code review, test generation, debugging, and refactoring.

\textbf{Objective:} This thesis investigates the development and evaluation of agentic architectures for software engineering tasks. The central questions are: (1)~How can robust agent strategies for SE workflows be modeled systematically? (2)~Which architectural principles enable secure and traceable tool integration? (3)~Which metrics allow a realistic assessment of agentic systems under practical constraints?

\textbf{Methodology:} Based on a structured literature review, a reference architecture is developed that integrates planning, tool usage, episodic memory, and safety mechanisms. A prototypical system is implemented in Python and evaluated using reproducible software-engineering scenarios.

\textbf{Results:} For the selected evaluation scenarios, the proposed architecture achieves a success rate of \pct{73} in automated refactoring and reduces mean processing time by \pct{45}. Reflection mechanisms improve robustness against faulty tool outputs by \pct{28}, while optimized context management reduces token usage by up to \pct{40}. These values should be interpreted as indicative prototype results rather than as directly comparable public leaderboard scores.

\textbf{Contribution:} The thesis provides a practical reference architecture, concrete implementation guidelines, and a multidimensional evaluation scheme for agentic SE systems. It identifies technical opportunities, names methodological limitations, and situates the findings within the current research landscape on agentic software engineering.%
\par
\endgroup

\cleardoublepage
